\documentclass[../../main.tex]{subfiles}

\begin{document}

\section{弱拓扑的度量化}\label{section:ch15-metrizability}

若 Banach 空间的闭单位球上的弱拓扑是可度量化的, 则 Eberlein-\v{S}mulian 定理是度量空间中的紧性与序列紧性的等价性的直接推论. 为更好地理解该定理, 我们现在证明弱拓扑的一些可度量化性质. 第一个定理给出分析学者不能仅停留在度量空间的好的理由.

\begin{theorem}\label{theorem:ch15-nonmetrizable}
令 $X$ 为无穷维赋范线性空间. 则 $X$ 上的弱拓扑与 $X^{*}$ 上的弱 $^{*}$ 拓扑都是不可度量化的.
\end{theorem}

\begin{proof}
为证明 $X$ 上的弱拓扑不可度量化, 我们用反证法. 否则, 存在度量 $\rho:X\times X\to[0,\infty]$ 诱导了 $X$ 上的弱拓扑. 固定自然数 $n$. 考虑 $0$ 的弱邻域 $\{x\in X\mid\rho(x,0)<1/n\}$. 我们可以选取 $X^{*}$ 的有限子集 $F_{n}$ 与 $\varepsilon_{n}>0$ 使得
\begin{align*}
\{x\in X\mid\text{对所有 }\psi\in F_{n},\ |\psi(x)|<\varepsilon_{n}\}\subseteq\{x\in X\mid\rho(x,0)<1/n\}.
\end{align*}
定义 $W_{n}$ 为 $F_{n}$ 的线性扩张. 则
\begin{align}
\bigcap_{\psi\in W_{n}}\ker\psi\subseteq\{x\in X\mid\rho(x,0)<1/n\}.\label{align:ch15-nonmet-1}
\end{align}
由于 $X$ 是无穷维的, 从 Hahn-Banach 定理得出 $X^{*}$ 也是无穷维的. 选取 $\psi_{n}\in X^{*}\sim W_{n}$. 我们从前一章的命题 4 推出, 存在 $x_{n}\in X$ 使得 $\psi_{n}(x_{n})\neq 0$, 而对所有 $\psi\in F_{n}$, $\psi(x_{n})=0$. 定义 $u_{n}=n\cdot x_{n}/\|x_{n}\|$. 观察到 $\|u_{n}\|=n$, 根据\eqref{align:ch15-nonmet-1}, $\rho(u_{n},0)<1/n$. 因此 $\{u_{n}\}$ 是 $X$ 中弱收敛到 $0$ 的无界序列. 这与前一章的定理 12 矛盾. 因此弱拓扑是不可度量化的.

为证明 $X^{*}$ 上的弱 $^{*}$ 拓扑是不可度量化的, 我们再次用反证法. 否则, 存在度量 $\rho^{*}:X^{*}\times X^{*}\to[0,\infty]$ 诱导了 $X^{*}$ 上的弱 $^{*}$ 拓扑. 固定自然数 $n$. 考虑 $0$ 的弱 $^{*}$ 邻域 $\{\psi\in X^{*}\mid\rho^{*}(\psi,0)<1/n\}$. 我们可以选取 $X$ 的有限子集 $A_{n}$ 与 $\varepsilon_{n}>0$ 使得
\begin{align*}
\{\psi\in X^{*}\mid\text{对所有 }x\in A_{n},\ |\psi(x)|<\varepsilon_{n}\}\subseteq\{\psi\in X^{*}\mid\rho^{*}(\psi,0)<1/n\}.
\end{align*}
定义 $X_{n}$ 为 $A_{n}$ 的线性扩张. 则
\begin{align}
\{\psi\in X^{*}\mid\text{对所有 }x\in X_{n},\psi(x)=0\}\subseteq\{\psi\in X^{*}\mid\rho^{*}(\psi,0)<1/n\}.\label{align:ch15-nonmet-2}
\end{align}
由于 $X_{n}$ 是有限维的, 它是闭的且是 $X$ 的真子空间, 这是因为 $X$ 是无穷维的. 从前一章的推论 11 得出, 存在在 $X_{n}$ 上等于零的非零泛函 $\psi_{n}\in X^{*}$. 定义 $\varphi_{n}=n\psi_{n}/\|\psi_{n}\|$. 观察到 $\|\varphi_{n}\|=n$, 且由\eqref{align:ch15-nonmet-2}, $\rho^{*}(\varphi_{n},0)<1/n$. 因此 $\{\varphi_{n}\}$ 是 $X^{*}$ 中逐点收敛到 $0$ 的无界序列. 这与一致有界定理矛盾. 因此 $X^{*}$ 上的弱 $^{*}$ 拓扑是不可度量化的.
\end{proof}

\begin{theorem}\label{theorem:ch15-W-metrizable}
令 $X$ 为赋范线性空间而 $W$ 是 $X^{*}$ 的分离 $X$ 中的点的可分子空间. $X$ 的闭单位球 $B$ 上的 $W$ 弱拓扑是可度量化的.
\end{theorem}

\begin{proof}
由于 $W$ 是可分的, $B^{*}\cap W$ 也是可分的, 其中 $B^{*}$ 是 $X^{*}$ 的闭单位球. 选取 $B^{*}\cap W$ 的可数稠密子集 $\{\psi_{k}\}_{k=1}^{\infty}$. 定义 $\rho:B\times B\to\mathbf{R}$ 为
\begin{align*}
\text{对所有 }u,\ v\in B,\quad\rho(u,v)=\sum_{k=1}^{\infty}\frac{1}{2^{k}}\cdot\left|\psi_{k}(u-v)\right|.
\end{align*}
由于每个 $\psi_{k}$ 属于 $B^{*}$, 这是适当定义的. 我们首先宣称 $\rho$ 是 $B$ 上的度量. $\rho$ 的对称性与三角不等式承袭自 $\psi_{k}$ 的线性. 另一方面, 由于 $W$ 分离 $X$ 中的点, $B^{*}\cap W$ 的任何稠密子集也分离 $X$ 中的点. 因此, 对 $u,\ v\in B$ 且 $u\neq v$, 存在自然数 $k$ 使得 $\psi_{k}(u-v)\neq 0$, 因此 $\rho(u,\ v)>0$. 因此 $\rho$ 是 $B$ 上的度量. 观察到对每个自然数 $n$, 由于每个 $\psi_{k}$ 属于 $B^{*}$,
\begin{align}
\text{对所有 }x\in B,\quad\frac{1}{2^{n}}\left[\sum_{k=1}^{n}\left|\psi_{k}(x)\right|\right]\leqslant\rho(x,0)\leqslant\sum_{k=1}^{\infty}\left|\psi_{k}(x)\right|+1/2^{n}.\label{align:ch15-W-metrizable}
\end{align}
我们把从这些不等式与 $\{\psi_{k}\}_{k=1}^{\infty}$ 在 $B^{*}\cap W$ 中的稠密性推出 $\{x\in B\mid\rho(x,\ 0)<1/n\}$ 是 $B$ 上的 $W$ 弱拓扑在原点的基留作练习. 因此由度量 $\rho$ 诱导的拓扑是 $B$ 上的 $W$ 弱拓扑.
\end{proof}

\begin{corollary}\label{corollary:ch15-metrizable}
令 $X$ 为赋范线性空间.
\begin{enumerate}[(i)]
\item 若 $X^{*}$ 是可分的, 则 $X$ 的闭单位球 $B$ 上的弱拓扑可度量化.
\item 若 $X$ 是可分的, 则 $X^{*}$ 的闭单位球 $B^{*}$ 上的弱 $^{*}$ 拓扑可度量化.
\end{enumerate}
\end{corollary}

\begin{theorem}\label{theorem:ch15-reflexive-metrizable}
令 $X$ 为自反 Banach 空间. 则闭单位球 $B$ 上的弱拓扑可度量化当且仅当 $X$ 是可分的.
\end{theorem}

\begin{proof}
由于 $X$ 是自反的, 前一章的定理 14 告诉我们若 $X$ 是可分的, $X^{*}$ 也如此. 因此, 根据前一个推论, 若 $X$ 是可分的, 则 $B$ 上的弱拓扑是可度量化的. 反过来, 假定 $B$ 上的弱拓扑是可度量化的. 令 $\rho:B\times B\to[0,\infty]$ 为 $B$ 上诱导弱拓扑的度量. 令 $n$ 为自然数. 我们可以选取 $X^{*}$ 的有限子集 $F_{n}$ 与 $\varepsilon_{n}>0$ 使得
\begin{align*}
\{x\in B\mid\text{对所有 }\psi\in F_{n},\ |\psi(x)|<\varepsilon_{n}\}\subseteq\{x\in B\mid\rho(x,0)<1/n\}.
\end{align*}
因此,
\begin{align}
\left[\bigcap_{\psi\in F_{n}}\ker\psi\right]\cap B\subseteq\{x\in B\mid\rho(x,0)<1/n\}.\label{align:ch15-reflexive-met}
\end{align}
定义 $Z$ 为 $\bigcup_{n=1}^{\infty}F_{n}$ 的闭线性扩张. 则 $Z$ 是可分的, 由于 $\bigcup_{n=1}^{\infty}F_{n}$ 中泛函的具有有理系数的有限线性组合是 $Z$ 的可数稠密子集. 我们宣称 $Z=X^{*}$. 否则, 前一章的推论 11 告诉我们存在非零 $S\in(X^{*})^{*}$, 它在 $Z$ 上消失. 由于 $X$ 是自反的, 存在某个 $x_{0}\in X$ 使得 $S=J(x_{0})$. 因此 $x_{0}\neq 0$ 且对所有 $\psi\in Z$, $\psi(x_{0})=0$. 根据\eqref{align:ch15-reflexive-met}, 对所有 $\rho(x_{0},\ 0)<1/n$, $x_{0}\neq 0$ 但 $\rho(x_{0},\ 0)=0$. 这是一个矛盾. 因此 $X^{*}$ 是可分的. 前一章的定理 13 告诉我们 $X$ 也是可分的.
\end{proof}

\begin{exercise}\label{exercise:ch15-16}
证明无穷维赋范线性空间的对偶也是无穷维的.
\end{exercise}

\begin{exercise}\label{exercise:ch15-17}
通过证明不等式\eqref{align:ch15-W-metrizable}蕴涵度量 $\rho$ 诱导 $W$ 弱拓扑, 完成定理\cref{theorem:ch15-W-metrizable}的证明的最后步骤.
\end{exercise}

\begin{exercise}\label{exercise:ch15-18}
令 $X$ 为 Banach 空间, $W$ 是它的对偶 $X^{*}$ 的闭子空间, 而 $\psi_{0}$ 属于 $X^{*}\sim W$. 证明: 若 $W$ 是有限维的或者 $X$ 是自反的, 则存在 $X$ 中的向量 $x_{0}$, 使得 $\psi_{0}(x_{0})\neq 0$ 但对所有 $\psi\in W$, $\psi(x_{0})=0$. 给出 $X^{*}$ 的一个无穷维闭子空间 $W$ 使得该分离性质不成立.
\end{exercise}

\end{document}